\documentclass{article}

\usepackage[sglblindworkshop]{neurips_2026}
\workshoptitle{Machine Learning for Systems}
\usepackage[T1]{fontenc}
\usepackage[utf8]{inputenc}
\usepackage{microtype}
\usepackage{amsmath}
\usepackage{amssymb}
\usepackage{booktabs}
\usepackage{multirow}
\usepackage{xcolor}
\usepackage[hidelinks]{hyperref}
\usepackage{url}

\newcommand{\system}{\textsc{FamilyPhase}}
\newcommand{\oracle}{\textsc{Oracle}}

\title{One Counter per Resource Family: Leakage-Aware Evaluation of Compact PMU Phase Prediction}

% Confirm the author list and institutional email before uploading to OpenReview.
\author{%
  Keshav Krishna \\
  Department of Computer Science \\
  New York University \\
  \texttt{keshaviitropar@gmail.com}
}

\begin{document}

\maketitle

\begin{abstract}
Machine learning can turn hardware performance counters into online workload-phase
predictors, but standard evaluation choices can make a predictor appear much more
deployable than it is. Randomly splitting correlated intervals leaks executions,
repeating the current offline phase creates an oracle baseline, overall accuracy hides
missed transitions, and model parameters omit counter discretization and history
state. We present \system, a compact and auditable interface in which a train-only
full-counter teacher discovers a global resource-pressure label while a deployable
student observes one validation-selected counter state per available resource family.
The evaluation groups complete co-running executions, reports transition-event
metrics and group-bootstrap uncertainty, distinguishes oracle from deployable inputs,
and budgets the complete online state. In an audited 60-run PARSEC pilot, a depth-six
tree over five 20-interval histories reaches 98.21\% held-out accuracy using an
estimated 201.25 bytes. Against a matched current-state-only tree, history improves
accuracy by 0.60 percentage points (paired 95\% CI: 0.29--0.92) and transition-event
F1 from 28.6\% to 43.9\%. The audit also exposes small transition counts, unstable
three-state clustering, and poor PMU enabled time. The main result is therefore an
evaluation protocol and reproducible study design, not an uncalibrated energy claim.
\end{abstract}

\section{Why phase-prediction accuracy is easy to overstate}

Hardware performance-monitoring units (PMUs) offer direct signals about cache,
memory, branch, and execution behavior. Learned phase predictors can use these signals
for scheduling, cache adaptation, or performance control
\citep{sherwood2003phase,perelman2006parallel,kim2017p4,lozano2022phase}. A practical
predictor, however, must work with few physical counter slots and a small online state.
It must also predict a label that is not already available at deployment.

Four common shortcuts obscure those requirements. First, a random interval split puts
near-identical samples from one execution in training and test. In multicore data it
can also separate two processes that ran together. Second, persistence and Markov
baselines are often given the offline teacher's true current phase. Such baselines
measure label persistence but cannot be implemented without a separate online phase
estimator. Third, stable intervals dominate: a model can achieve high accuracy while
never identifying a phase change. Finally, reporting only tree or network parameters
omits the state needed to discretize counters and store history.

We study the following deliberately constrained question:

\begin{quote}
Can one counter state from each resource family predict the next global phase, and can
that claim survive execution-disjoint evaluation, deployable baselines, transition
metrics, PMU-quality auditing, and a complete state budget?
\end{quote}

Our contribution is methodological. We combine (1) a train-only teacher/student
interface with validation-only counter selection, (2) leakage-aware splitting and
execution-group bootstrap inference, and (3) an information-set audit that prevents
oracle phase labels from being presented as online inputs. An accompanying collector
enumerates all workload pairs, records exact event encodings and enabled time, and
emits the tables needed to replace the preliminary evidence below.

\section{Relation to prior phase learning}

Program-phase methods traditionally summarize basic-block execution and cluster the
resulting vectors offline \citep{sherwood2001bbda,sherwood2002simpoint}. Subsequent
work compared phase representations, predicted recurring phases, and extended phase
analysis to parallel programs
\citep{dhodapkar2003comparing,sherwood2003phase,perelman2006parallel}. More recent
systems use hardware counters and learned models to make prediction inexpensive
enough for online control \citep{nomani2015predicting,kim2017p4,lozano2022phase}.
\system{} is complementary: its novelty is not another high-capacity predictor, but
an evaluation boundary around a deliberately small one-counter-per-family student.
That boundary matters because physical event slots are scarce, multiplexed events can
distort labels, and discretization and history consume state even when model
parameters are tiny \citep{eyerman2006cpi}. We therefore ask whether the compact
interface survives execution-disjoint and workload-disjoint tests before connecting
it to a resource controller.

\section{A compact, explicit information interface}

\paragraph{Global teacher.}
For run $r$ and interval $t$, let $x_{r,t}\in\mathbb{R}^{d}$ contain the safe
nonnegative PMU counts available on the host. Timing-derived quantities such as IPC,
CPI, elapsed time, and requested frequency are excluded so that a future controller
does not feed its own action back into the label. Training medians fill missing
values; $\log(1+x)$ and training-only standardization limit count skew. The offline
teacher fits $k$-means on training rows,
\begin{equation}
  c_{r,t}=\arg\min_j\lVert \widetilde{x}_{r,t}-\mu_j\rVert_2^2,
\end{equation}
and freezes its preprocessing and centroids for validation and test. Cluster IDs are
ordered by a resource-pressure score for interpretation. They are descriptive regions
of the observed counter space, not ground-truth bottleneck or power labels.

\paragraph{Reduced student input.}
Candidate events are assigned to L1, L2, LLC, branch/control, core/FP, and
memory/off-core families. A depth-limited singleton tree is trained for each supported
event, and validation accuracy plus high-pressure recall selects one event
$e_f^\star$ per family. Test rows never select counters. A train-only one-dimensional
three-cluster model converts each selected count into
$s_{f,t}\in\{0,1,2\}$. The deployable history model predicts
\begin{equation}
 h([s_{1,t-H+1:t},\ldots,s_{F,t-H+1:t}])\rightarrow c_{r,t+1}
\end{equation}
with a shallow CART tree \citep{breiman1984cart}. All family streams share the same
global target; a family name identifies an input, not a separate phase ground truth.

\paragraph{Information-set audit.}
Table~\ref{tab:inputs} separates diagnostic and deployable models. Persistence,
state-conditioned majority, Markov, run-length Markov, and duration models consume
the true offline $c_{r,t}$ and are labeled oracle diagnostics. The fair history
ablation is a tree with identical depth and leaf constraints that observes only the
current vector $s_{1..F,t}$. A training-prior majority model is a second deployable
sanity check.

\begin{table}[t]
  \caption{Inputs available at prediction time. ``Teacher phase'' means the offline
  full-counter label and is not a deployable input.}
  \label{tab:inputs}
  \centering
  \small
  \begin{tabular}{lll}
    \toprule
    Model & Prediction-time input & Deployable? \\
    \midrule
    Majority & training label prior & yes \\
    Persistence/Markov & true teacher $c_{r,t}$ & no: oracle \\
    Current-state tree & selected $s_{1..F,t}$ & yes \\
    History tree & selected $s_{1..F,t-H+1:t}$ & yes \\
    \bottomrule
  \end{tabular}
\end{table}

\paragraph{Complete analytical state budget.}
For a tree with $I$ internal nodes, $L$ leaves, $D$ inputs, and $N=I+L$ nodes, we
estimate
\begin{equation}
 B_{\mathrm{tree}}=
 \frac{I(\lceil\log_2D\rceil+16+2\lceil\log_2N\rceil)+2L}{8}.
\end{equation}
The online total adds $\lceil2FH/8\rceil$ bytes for two-bit history and $4F$ bytes for
two 16-bit discretizer thresholds per family. This is storage, not synthesized area or
energy.

\section{Evaluation protocol}

\paragraph{Leakage control.}
The split unit is a complete execution scenario. All intervals and repetitions from a
configuration remain together; for co-runners, every process in the concurrent group
receives the same split. Leave-one-workload-out evaluation moves every scenario
containing the held-out workload to test, including its partners. Counter selection
uses validation data only. Preprocessing, teacher centroids, counter discretizers, and
student models use training data only.

\paragraph{Metrics and uncertainty.}
We report overall accuracy, macro F1, balanced accuracy, high-pressure recall, and
separate stable- and transition-destination accuracy. Treating a predicted label
change as an event gives transition precision, recall, F1, and false-alarm rate. The
95\% intervals resample complete execution/concurrent-group IDs rather than intervals.
History versus current-state comparisons use paired resampling of the same groups.
Clustering sensitivity covers $k\in\{2,3,4,5\}$ and five seeds using adjusted Rand
index (ARI), train distance, centroid separation, and minimum cluster share.

\paragraph{Audited pilot.}
The reproducible pilot contains 60 executions of five PARSEC workloads
\citep{bienia2008parsec}, four thread counts, and three repetitions per configuration.
It was collected at 10~ms on a dual-socket Intel Xeon E5-2680~v2 host using Linux
\texttt{perf} 5.14. The merged table contains 9,409 intervals and eight usable
counters; L2 was unavailable, leaving five online families. The primary grouped split
contains 5,642 training and 1,341 test history windows, but only nine test execution
groups and 31 true transitions. We therefore call every number below preliminary.

\section{What the pilot establishes---and what it does not}

Table~\ref{tab:pilot} shows why the information audit matters. Oracle persistence
reaches 97.69\% accuracy because the target rarely changes, yet it correctly predicts
no transition destination. The deployable current-state tree reaches 97.61\% accuracy
and 28.6\% transition-event F1. Adding 20 intervals of history reaches 98.21\%
accuracy and 43.9\% event F1. The paired accuracy improvement is 0.60 points (95\% CI:
0.29--0.92). Transition-destination accuracy improves by 12.95 points, but its
interval is 0.0--26.92; the pilot is underpowered for a strong transition claim.

\begin{table}[t]
  \caption{Pilot configuration-group holdout. Accuracy CIs resample nine test
  executions. Transition F1 treats a label change as an event.}
  \label{tab:pilot}
  \centering
  \small
  \setlength{\tabcolsep}{4pt}
  \begin{tabular}{lrrrr}
    \toprule
    Model & Accuracy (95\% CI) & Macro F1 & Trans. acc. & Trans. F1 \\
    \midrule
    Majority & 77.03 (70.77--82.69) & 29.01 & 48.4 & 9.4 \\
    Oracle persistence & 97.69 (97.20--98.14) & 95.97 & 0.0 & 0.0 \\
    Current-state tree & 97.61 (96.84--98.30) & 95.45 & 12.9 & 28.6 \\
    History tree & \textbf{98.21} (97.64--98.71) & \textbf{96.88} & \textbf{25.8} & \textbf{43.9} \\
    \bottomrule
  \end{tabular}
\end{table}

Workload-disjoint evaluation is harder. Averaged descriptively over five
leave-one-workload-out folds, the history tree reaches 94.60\% accuracy and 43.9\%
transition-destination accuracy, versus 93.01\% and 22.5\% for the current-state tree.
The bodytrack holdout falls to 82.24\% history accuracy, demonstrating that the
configuration split alone is insufficient. Accuracy improves reliably in four folds;
canneal's 0.07-point difference has a paired 95\% CI of $-0.38$--0.63.

\begin{table}[t]
  \caption{Pilot leave-one-workload-out results (\%). Transition accuracy is
  destination accuracy on true changes. Means are descriptive across five folds.}
  \label{tab:loo}
  \centering
  \small
  \setlength{\tabcolsep}{4pt}
  \begin{tabular}{lrrrr}
    \toprule
    Held out & \multicolumn{2}{c}{Accuracy} & \multicolumn{2}{c}{Trans. accuracy} \\
    \cmidrule(lr){2-3}\cmidrule(lr){4-5}
    & Current & History & Current & History \\
    \midrule
    blackscholes & 97.68 & 98.24 & 21.43 & 40.48 \\
    bodytrack    & 76.47 & 82.24 & 21.04 & 47.26 \\
    canneal      & 99.06 & 98.99 & 17.65 & 35.29 \\
    fluidanimate & 98.38 & 99.15 & 12.50 & 37.50 \\
    freqmine     & 93.45 & 94.40 & 40.00 & 59.09 \\
    \midrule
    Mean         & 93.01 & 94.60 & 22.52 & 43.93 \\
    \bottomrule
  \end{tabular}
\end{table}

Three audit findings prevent a broader claim. First, $k=2$ is almost invariant across
seeds (mean ARI 0.9997), while $k=3$ has mean ARI 0.927 and a smallest cluster of only
2.1\% of training rows. Low/moderate/high semantics do not by themselves justify the
less stable choice. Second, the eight teacher events average 90.6--91.0\% enabled
time, but their fifth percentiles are only 64--66\%, and roughly one-third of samples
fall below 90\%. Multiplexing quality must accompany model accuracy. Third, 31 test
transitions cannot support an end-to-end control conclusion.

The complete online-state estimate is 64.88 bytes for the current-state tree and
201.25 bytes for the history tree. The latter includes 25 bytes of history and 20
bytes of thresholds. These values show compactness but not RTL area, latency, or
energy. We intentionally remove an earlier proxy that mapped phase labels to arbitrary
cache-power factors: no cache capacity, miss latency, energy, or slowdown model had
been calibrated.

\section{Real-data study and ML-for-systems lessons}

The accompanying one-command study targets three regimes: a single multithreaded
process, two single-threaded co-runners, and two four-thread co-runners. Seven PARSEC
workloads produce all $\binom{7}{2}=21$ cross-workload pairs, with ten repetitions,
fixed CPU affinity, randomized execution-group order, exact platform/event manifests,
and PMU enabled-time summaries. It runs configuration-group and leave-one-workload-out
evaluation together, performs $k$/seed sensitivity, and emits paired current-versus-
history bootstraps. Large data, environments, and caches are confined to scratch
storage.

This design yields three lessons for ML-for-systems evaluation. \textbf{Information
sets are part of the model}: a simple Markov method with an unavailable state is not a
deployable baseline. \textbf{Correlation defines the split unit}: systems traces are
not independent examples, and co-runners cannot be separated. \textbf{Measurement
quality is model quality}: an offline teacher trained on poorly scheduled events can
produce a precise student of an unstable target. These checks precede architecture
search; a larger neural model cannot repair leakage or unavailable inputs.

The planned evidence does not guarantee cross-platform portability. PMU semantics,
event availability, workload mix, NUMA placement, and operating-system activity remain
host dependent. A later archival paper should add a second contemporary platform and,
if it claims control benefit, an actual actuator such as cache allocation or DVFS with
measured performance and energy. The present claim is intentionally narrower: a
compact phase interface and a protocol for evaluating it without granting the learner
information it will not have online.

\section{Conclusion}

Compact PMU histories can predict a persistent offline phase label, but accuracy alone
does not establish a useful online system. \system{} makes the target, prediction-time
inputs, split unit, transition behavior, PMU quality, and full state budget explicit.
The pilot finds a small paired benefit from history and, more importantly, reveals the
conditions under which that result may fail. We release the corrected study workflow
so the final claim can be decided by workload-disjoint multicore evidence rather than
interval leakage or an oracle baseline.

\clearpage
\bibliographystyle{plainnat}
\bibliography{references}

\end{document}
